% element: 10-s033:e04
\BeleskeID{10-s033}
Ulazna mreža sa dodatnim otpornikom ka negativnom napajanju može da se zameni Tevenenovim ekvivalentom. Za ekvivalentni napon važi
\[E_T=V_I\frac{R_N}{R_B+R_N}-V_N\frac{R_B}{R_B+R_N}.\]
Da bi tranzistor dostigao granicu provođenja, taj napon treba da bude $V_{\gamma T}$. Zato za ulazni napon pre Tevenenove transformacije na pragu pišemo
\[V_{\gamma T}=V_{IL}\frac{R_N}{R_B+R_N}-V_N\frac{R_B}{R_B+R_N}.\]
Rešavanjem po $V_{IL}$:
\[V_{IL}=\frac{R_B+R_N}{R_N}\left(V_{\gamma T}+V_N\frac{R_B}{R_B+R_N}\right).\]
Raspodela otpornosti i negativno napajanje pomeraju ulazni prag; efektivni bazni otpornik je njihova paralelna veza.
